%% Paper 1 — elsarticle SKELETON (Computers and Electronics in Agriculture)
%% Structural skeleton only. Prose to be finalized from docs/WRITING_DOSSIER.md.
%% All numbers trace to docs/numbers_ledger.md. Citations are placeholders; verify
%% every \cite key against refs/references.bib (each marked VERIFY). No invented DOIs.
\documentclass[preprint,12pt]{elsarticle}
\usepackage{amsmath,graphicx,booktabs,siunitx}
\journal{Computers and Electronics in Agriculture}

\begin{document}
\begin{frontmatter}
\title{Spatial skill is not temporal skill: an honest cross-validation audit of
satellite-driven wheat and sunflower yield prediction in Trakya, T\"urkiye}

\author[a]{Melih Kalkan}
\author[a]{\c{C}avdaro\u{g}lu} %% VERIFY full name/initials/affiliation
\address[a]{VERIFY: institution, Trakya, T\"urkiye}

\begin{abstract}
%% From manuscript/paper1.md Abstract (numbers verified in docs/numbers_ledger.md).
Satellite- and machine-learning-based crop-yield models are frequently reported with high
accuracy, but the cross-validation design often rewards spatial interpolation rather than the
temporal extrapolation operational use requires. Re-analysing winter wheat and sunflower over 29
Trakya districts (2004--2025) under a frozen, recompute-faithful protocol, we contrast
leave-one-year-out (LOYO, temporal), leave-one-district-out (LOILO, spatial), and spatiotemporal
cross-validation against climatology and persistence baselines. Under temporal validation no
machine-learning configuration beats climatology for wheat (best skill score $-0.172$, 95\,\% CI
$[-0.238,-0.109]$), whereas a multimodal model does for sunflower ($R^2=0.386$; skill score
$+0.224\,[+0.167,+0.276]$). We quantify a large, significant spatial-vs-temporal generalization gap
(e.g.\ wheat climate tier $\Delta R^2=+0.639\,[+0.545,+0.735]$), show NDVI's marginal value is
crop-specific via a matched-sample ablation, and document a senescence collapse in forward NDVI
forecasting. Cross-validated accuracy is not operational skill.
\end{abstract}

\begin{keyword}
crop yield prediction \sep spatial vs temporal cross-validation \sep generalization gap \sep
NDVI \sep climatology baseline \sep operational validation
\end{keyword}
\end{frontmatter}

%% Highlights (submit separately; 3-5 bullets, <=85 chars) -- see manuscript/paper1.md.

\section{Introduction}\label{sec:intro}
% paper1.md Sec 1. Cite \cite{vanklompenburg2020review,roberts2017crossval,ploton2020spatial}. [VERIFY keys]

\section{Study area and data}\label{sec:data}
% paper1.md Sec 2. TUIK yields; Sentinel-2 \cite{drusch2012sentinel2}; GEE \cite{gorelick2017gee};
% SoilGrids \cite{poggio2021soilgrids}. Tiers A/B/C; n=589/576 (climate), 213/209 (NDVI). [VERIFY]

\section{Methods}\label{sec:methods}
% paper1.md Sec 3 + docs/methods_as_implemented.md. CV regimes; models; baselines B0-B3;
% bootstrap (iid+cluster); paired Wilcoxon + rank-biserial; matched ablation; Moran's I
% \cite{moran1950notes}; FLOV. Reproduction fidelity max|delta|=0.0.

\section{Results}\label{sec:results}
% paper1.md Sec 4. Tables 1-2, Figs 1-6.
\begin{table}[t]\centering\caption{Cross-validated performance (excerpt; full in T1\_master\_results.csv).}
\label{tab:t1}\begin{tabular}{llrrr}\toprule
Crop & Model (LOYO) & $R^2$ & RMSE & SS vs B0\\\midrule
Wheat & B0 climatology & 0.213 & 61.7 & 0.000\\
Wheat & Layer C XGBoost & -0.311 & 80.7 & -0.172\\
Sunflower & B0 climatology & 0.033 & 50.0 & 0.000\\
Sunflower & Layer C GPR & 0.386 & 42.3 & +0.225\\\bottomrule
\end{tabular}\end{table}

\section{Discussion}\label{sec:discussion}
% paper1.md Sec 5.

\section{Limitations}\label{sec:limitations}
% paper1.md Sec 6. EVR_01 now REAL surveyed coords (0.62 ha, single small field, subpixel);
% small-n LOYO; single region; EVR_02-05 still placeholder.

\section{Conclusion}\label{sec:conclusion}
% paper1.md Sec 7.

\section*{Data and code availability}
% paper1_generalization/ ; run_all.py.

\section*{Acknowledgements}
T\"UB\.ITAK 2209-A %% [VERIFY grant id]

\bibliographystyle{elsarticle-num}
\bibliography{../refs/references}
\end{document}
